\documentclass[11pt,a4paper]{ltjsarticle}
\usepackage{luatexja-fontspec}
\usepackage{geometry}
\usepackage{booktabs}
\usepackage{multirow}
\usepackage{enumitem}
\usepackage{amsmath}
\usepackage{fancyhdr}

\geometry{margin=22mm}
\setmainjfont[BoldFont=Hiragino Kaku Gothic ProN]{Hiragino Mincho ProN}
\setsansjfont{Hiragino Kaku Gothic ProN}

\pagestyle{fancy}
\fancyhf{}
\lhead{修士論文 骨子（和文版）}
\rhead{2026 年 7 月 13 日}
\cfoot{\thepage}
\renewcommand{\headrulewidth}{0.2pt}

\newcommand{\rk}{\textrm{Recall@}K}
\setlist{nosep, leftmargin=1.6em}
\setlength{\parskip}{0.2\baselineskip}

\begin{document}

\begin{center}
{\sffamily\large 修士論文 骨子：文献から抽出した数式と文脈を用いた\\
物理モデル自動構築のための数式集合検索}\\
\vspace{0.2em}
{\small (Equation-Set Retrieval Using Equations and Context Extracted from the Literature for Automated Physical Model Building)}
\end{center}
\vspace{0.1em}
\hfill 作成者：M2\ 一洋
\vspace{0.4em}

\noindent
【本資料について】英語骨子（\texttt{thesis/master\_thesis/main.pdf}、16 ページ）のレビュー用和文要約である。章立て・段落構成は英語版と同一で、各箇条書きが本文の 1 段落（先頭文）に対応する。

%======================================================================
\section*{論文の主張（要約）}
%======================================================================
本論文は、「どんな物理モデルを作りたいか」（説明文と入出力変数）を入力すると、複数文献から集めた数式データベース（11{,}146 式・361 ソース）の中から、そのモデルを構成する\textbf{数式の集合}を検索して提示する手法を提案する。核心は、候補の式を 1 本ずつでなく\textbf{すでに選んだ集合との関係}で評価する集合特徴量である。$\rk$ は古典的検索の 0.521 から \textbf{0.743}（設定 A）へ、最難の DAE のみでも 0.466 から \textbf{0.689} へ、いずれも +0.22 改善する（Wilcoxon 符号順位検定 $p=0.00195$）。LLM 直接生成（0.23--0.30）に対しても優位性を示す。さらに、(i) 上乗せの源泉が構造的な集合特徴量であることの実証、(ii) 学習版 greedy による最難端の回復、(iii) 正解式数 $K$ を知らずに「解ける系」を返す自由度停止、までを一貫した流れとして示す。

%======================================================================
\section*{第 1 章　序論（Introduction）}
%======================================================================
\begin{enumerate}
  \item 物理モデルはデジタルツインの実現やプロセスの設計・運転に不可欠だが、構築には膨大な文献調査と試行錯誤を要する。
  \item 我々のグループはこれを自動化する AutoPMoB の開発を目指しており、変数同義性判定（ProcessBERT）などの要素技術を積み上げてきた。本研究はその次の段階＝\textbf{抽出済みの数式をモデルへ組み上げる}工程を担う。
  \item 物理モデルは変数を共有し連立して初めて解ける系（方程式数＝未知数の数、自由度ゼロ）なので、文章類似度だけの検索では足りない。
  \item そこで 2 段階検索（第 1 段 7 特徴量 → 第 2 段 集合特徴量 3 個＋MLP）に、学習版 greedy と自由度停止を加えた手法を提案する。
  \item 評価は 11{,}146 式 DB・解けることを保証した DAE ケース・ソース単位の層化分割（10 乱数）で行う。
  \item 貢献 3 点：(1) 解けることを構成的に保証したデータセット、(2) 集合を考慮した 2 段階検索（+0.22、古典 IR と LLM 直接生成の双方に優位）、(3) 機構の実証 → 学習版 greedy → 自由度停止。
\end{enumerate}

%======================================================================
\section*{第 2 章　関連研究（Related Work）}
%======================================================================
4 系統（AutoPMoB 要素技術／情報検索と再順位付け／LLM による知識の直接生成／数式検索）を概観し、各系統の末尾で本研究との相違＝\textbf{独立に関連する項目でなく、連立して解ける集合を返す}点に毎回戻る。実文献の引用は本文化の段階で検証のうえ追加する。

%======================================================================
\section*{第 3 章　手法（Set-aware Equation-Set Retrieval）}
%======================================================================
\begin{enumerate}
  \item \textbf{問題定式化}：入力 $q=$（説明文, 入力変数, 出力変数）、出力＝DB 中の数式の順位付き集合。自由度 $\mathrm{DoF}(S)=|\text{未知数}|-|S|$ が 0 のとき系は閉じる（解ける）。
  \item \textbf{第 1 段}：7 特徴量（説明文との TF-IDF 類似度・入出力変数との Jaccard・SVD 潜在類似度・入力被覆率・出力被覆率・変数適合率・ドメイン一致）で候補を絞る。これが比較対象の baseline でもある。
  \item \textbf{第 2 段}：集合特徴量 3 個 --- 補完性 gComp（未充足の入出力変数をどれだけ補うか）・一貫性 gCoh（既選択集合と変数をどれだけ共有するか）・ドメイン一致 gDom --- を加えた 10 特徴量を MLP で評点（reranker-10S）。
  \item \textbf{逐次選択 3 方式}：静的／推論時 greedy（選ぶたびに集合特徴量を再計算）／学習版 greedy（teacher forcing で訓練時も逐次文脈を与え、訓練と推論の条件を一致させる）。
  \item \textbf{自由度停止}：$K$ を与えずに、系が閉じた時点で自動停止する。
\end{enumerate}

%======================================================================
\section*{第 4 章　データセット（Equation Database and Test Case Construction）}
%======================================================================
\begin{enumerate}
  \item \textbf{数式 DB}：論文 329 報を LLM（Gemini 2.5 Pro）で抽出＋32 分野の基礎式 670 式を生成し、11{,}146 式・361 ソースに拡大（3{,}244 式から）。
  \item \textbf{ケース 2{,}838 件}：オリジナル 92・複数文献 731（規則ベース、LLM 不使用）・DAE 1{,}000・データ拡張 1{,}015（訓練のみ、評価から除外）。複数式ケースが 74.6\%。
  \item \textbf{DAE ケース生成}：加藤先生にご指摘いただいた手順（ODE を種に、変数を共有する式を追加し、方程式数＝未知数の正方系のみ採用。$X=1..10$ × 各 100 件）。説明文の生成のみ LLM（Claude）を使用。
  \item \textbf{層化分割}：ソース単位で 6:2:2 に分け（言い換えによる答え漏れを防止）、400 候補から 4 分布（数式数・入出力変数数・横断ソース数）の $L_1$ 距離最小の分割を選択。10 乱数で反復。
\end{enumerate}

%======================================================================
\section*{第 5 章　実験（Experiments）}
%======================================================================
設定 A（原型＋複数文献＋DAE、1{,}823 件）と設定 B（DAE のみ 1{,}000 件）。拡張込みで評価すると baseline が +0.101 水増しされ本手法の効果を過小評価するため、主要評価から除外する。比較 3 手法＝baseline（古典 IR）／reranker-10S（本手法）／LLM 直接生成（生成式の等価性を別 LLM で判定、各設定 $n\approx50$）。指標は $\rk$・MAP・Recall@20、自由度停止には完全一致率・集合 F1・閉包率。10 乱数の対応あり Wilcoxon 符号順位検定を行う（$n=10$ の両側最小 $p$ は 0.00195）。

%======================================================================
\section*{第 6 章　結果と考察（Results and Discussion）}
%======================================================================
流れ：\textbf{性能 → 機構 → 介入 → 自己停止 → 留意点}。
\begin{enumerate}
  \item \textbf{全体}（表 1）：設定 A で 0.521$\to$0.743、設定 B で 0.466$\to$0.689（いずれも $p=0.00195$、10 乱数全勝）。LLM 直接生成は 0.23--0.30 にとどまる一方、順位を問わない coverage は 0.34--0.57 ＝「生成はできるが上位 $K$ 件に並べられない」。
  \item \textbf{特性別}：数式数・入力変数数が多い難しいケースほど優位性が拡大（+0.03 $\to$ +0.21〜0.29）。ソース数は数式数と交絡しており単独の効果ではない。
  \item \textbf{機構}：補完性 +0.052・一貫性 +0.050 は各単独で有意、ドメイン一致は寄与ゼロ（$p=0.98$）。上乗せは閾値＋こぶ型（$X=3$ で +0.12）で、最難端では飽和する。
  \item \textbf{介入}：学習版 greedy が飽和を回復。DAE で静的 0.724 $<$ 推論 greedy 0.756 $<$ 学習版 0.765（学習版 vs 推論 $p=0.037$）、$X\ge8$ で静的比 +0.056。設定 A でも 0.757 $<$ 0.787 $<$ 0.797（greedy vs 静的は有意 $p=0.00195$、学習版 vs 推論の +0.010 は非有意 $p=0.13$。greedy 系は top-$k$=200 のため表 1 の値とは直接比較しない）。
  \item \textbf{自由度停止}：完全一致率は $K$ 既知の oracle と厳密に同一（DAE 0.368／設定 A 0.526）。閉包率は oracle を上回り（0.883 vs 0.781、0.936 vs 0.865）＝\textbf{解ける系を返す}保証。集合 F1 のコストは僅少（設定 A $-0.008$）。ただし本データでは $K\approx$ 出力変数数のため $K$ 予測自体は自明であり、主張は「自己停止＋閉包保証」に置く。
  \item \textbf{留意点}：DAE の説明文は LLM 生成、ソース数の交絡、分割方式は 400 候補から選んだ 1 つ、DB 内評価（DB にない式は検索できない）。
\end{enumerate}

\begin{table}[h]\centering
\small
\caption{全体結果（baseline・本手法は層化分割・10 乱数の平均、$\pm$ は乱数間の標準偏差）。LLM 直接生成は各設定の一様サンプル（A: $n=48$、B: $n=50$）の単発評価で、生成式の等価性を別の LLM（claude-opus-4-8）で判定して採点した。順位を問わない coverage は 0.419（A）・0.343（B）。}
\begin{tabular}{llccc}
\toprule
設定 & 手法 & $\rk$ & MAP & Recall@20 \\
\midrule
\multirow{3}{*}{A（1{,}823 件）}
 & baseline（古典 IR） & $0.521\pm0.045$ & 0.581 & 0.836 \\
 & LLM 直接生成（外部比較） & $0.254$ & --- & --- \\
 & \textbf{reranker-10S（本手法）} & $\mathbf{0.743\pm0.054}$ & \textbf{0.863} & \textbf{0.909} \\
\midrule
\multirow{3}{*}{B（DAE のみ 1{,}000 件）}
 & baseline（古典 IR） & $0.466\pm0.037$ & 0.512 & 0.721 \\
 & LLM 直接生成（外部比較） & $0.226$ & --- & --- \\
 & \textbf{reranker-10S（本手法）} & $\mathbf{0.689\pm0.033}$ & \textbf{0.832} & \textbf{0.838} \\
\bottomrule
\end{tabular}
\end{table}

%======================================================================
\section*{第 7 章　結論（Conclusion）}
%======================================================================
貢献 3 点（データセット／集合を考慮した 2 段階検索／機構実証と自己停止）を要約し、今後の課題として scheduled sampling による学習版 greedy の改善、DB 外の式を含む実モデリング課題での検証、AutoPMoB パイプラインへの統合を挙げる。

\end{document}
